\chapter{Glossary}
\label{chap:Glossary}

The following glossary defines the principal technical, legal, and procedural terms used throughout \textit{Prudentia Aedificatoria}. Terms are arranged alphabetically. Where a term has both a general meaning and a specific legal or regulatory meaning, the regulatory meaning is given precedence. Cross-references to relevant chapters are provided in square brackets.

\begin{description}

\item[ACM (Asbestos-Containing Material)] Any material or product that contains asbestos of any type (chrysotile, amosite, crocidolite, or others). ACMs must be identified by a survey, their condition assessed, and a management plan produced. Work disturbing ACMs above a certain threshold requires a licensed asbestos contractor. [Ch.12]

\item[Appointed Person (Lifting Operations)] The competent person appointed under LOLER 1998 to take overall responsibility for the planning and supervision of lifting operations on a construction site. The appointed person must have sufficient knowledge, skill, and experience to plan lifts, select equipment, and oversee safe execution. [Ch.15]

\item[Asbestos Analyst] A person holding a relevant accreditation (typically UKAS-accredited to ISO 17025) for the collection and analysis of asbestos samples, including four-stage clearance air testing following licensed asbestos removal work. [Ch.12]

\item[Banksman] A trained operative who guides a crane, excavator, or other large item of mobile plant by agreed hand signals or radio communication, ensuring the movement of the machine does not endanger other persons. Banksmen must be specifically trained and must not be assigned other duties while acting in the banksman role. [Ch.6, Ch.15]

\item[Bow-Tie Model] A risk assessment and incident investigation tool that maps the hazard, the top event (the undesired incident), the threat pathways (left side, with prevention barriers) and the consequence pathways (right side, with mitigation barriers). Widely used in high-hazard industries and increasingly applied in construction incident investigation. [Ch.29]

\item[CAT (Cable Avoidance Tool)] An electronic instrument used to detect the electromagnetic field produced by buried electrical cables and pipes carrying a current. Used in conjunction with a signal generator (Genny) to identify the position of buried services before excavation. CAT surveys must be conducted by a trained and competent operative; results must be marked and recorded. [Ch.6, Ch.22]

\item[CDM 2015] The Construction (Design and Management) Regulations 2015. The principal UK regulatory framework for health and safety in construction, setting out the duties of the client, principal designer, principal contractor, contractor, designer, and workers. [Ch.27]

\item[CHAS (Contractors Health and Safety Assessment Scheme)] A UK SSIP-accredited pre-qualification scheme that assesses contractors against the SSIP core criteria set. CHAS accreditation is widely accepted by principal contractors and public sector clients as evidence of basic health and safety competence. [Ch.28]

\item[CITB (Construction Industry Training Board)] The industry training board for the construction sector in Great Britain, responsible for funding and supporting construction training and apprenticeships. CITB administers the CSCS card scheme in partnership with industry. [Ch.27]

\item[Client (CDM)] The organisation or individual for whom a construction project is carried out. Under CDM 2015, the client has specific duties to ensure suitable arrangements for managing the project, appoint Principal Designer and Principal Contractor, provide pre-construction information, and ensure the Health and Safety File is prepared and maintained. [Ch.27]

\item[CLOCS (Construction Logistics and Cyclist Safety)] An industry standard for construction logistics that requires principal contractors and their supply chain logistics providers to manage the safety of vulnerable road users (cyclists, pedestrians, motorcyclists) who interact with delivery vehicles. [Ch.19, Ch.28]

\item[Competent Person] A person who has sufficient training, experience, knowledge, and other qualities to perform a specific function safely. The specific competence requirements vary by function; for CDM purposes, the concept of competence covers knowledge, skills, experience, and organisational capability. [Throughout]

\item[Confined Space] A place that is substantially enclosed (though not always entirely), and where there is a reasonably foreseeable risk of serious injury from hazardous substances or conditions within the space (or nearby). Examples on construction sites include manholes, sewers, inspection chambers, tanks, tunnels, and enclosed plant rooms. Defined by the Confined Spaces Regulations 1997. [Ch.16]

\item[Consequence] The severity of harm that would result if a hazard causes an injury or illness. In the 5$\times$5 risk matrix used throughout this book, consequence is rated 1 (minor, first aid) to 5 (catastrophic, multiple fatalities). [Ch.3]

\item[Construction Phase Plan (CPP)] A document prepared by the Principal Contractor (or sole contractor) before the construction phase starts, describing the health and safety management arrangements for the project. The CPP must be developed and maintained throughout the construction phase. Required by CDM 2015. [Ch.27]

\item[COSHH (Control of Substances Hazardous to Health Regulations 2002)] The regulations that require employers to prevent or adequately control the exposure of employees and others to substances hazardous to health. COSHH assessments must be produced for all hazardous substances used in or created by the work. [Ch.11]

\item[CPCS (Construction Plant Competence Scheme)] An accreditation scheme for construction plant operators, assessing the competence of crane operators, excavator operators, MEWP operators, and others against industry standards. CPCS cards are recognised evidence of plant operator competence. [Ch.9, Ch.15]

\item[CSCS (Construction Skills Certification Scheme)] The industry scheme providing cards to workers and managers as evidence of health and safety training and qualifications appropriate to their occupation. Card levels include: Green (Labourer), Blue (Skilled Worker), Gold (Supervisor), Black (Manager). [Ch.27]

\item[CTE (Cut-to-Length Excavator)] An excavator configured to cut and process tree-length timber, commonly used in forestry and infrastructure clearance. Referenced for completeness where vegetation clearance is part of the pre-construction phase. [Ch.6]

\item[Dangerous Occurrence] One of 27 specified categories of near-miss events defined in Schedule 2 to RIDDOR 2013 that must be reported to the relevant enforcing authority even where no injury has occurred. Examples relevant to construction include: scaffold collapse; crane overturning; unintended collapse of a building or structure; unintended release of flammable gas or highly flammable liquid. [Ch.29]

\item[Designer (CDM)] Any organisation or individual who, as part of a business, prepares or modifies a design or arranges for or instructs any person under their control to do so. Designers under CDM include architects, structural engineers, building services engineers, interior designers, and temporary works designers. Designers have specific duties to eliminate, reduce, or control foreseeable risks arising from their designs. [Ch.27]

\item[Duty Holder] Any person or organisation that has specific health and safety duties under CDM 2015 or other health and safety legislation. Under CDM 2015, duty holders include the client, principal designer, principal contractor, contractor, designer, and worker. [Ch.27]

\item[EHO (Environmental Health Officer)] An officer of a Local Authority with enforcement powers under health and safety legislation for non-CDM-notifiable construction and for fixed premises. EHOs have the same inspection and enforcement powers as HSE inspectors for premises within their jurisdiction. [Ch.2]

\item[EAP (Employee Assistance Programme)] A confidential support service providing employees with access to counselling, financial advice, legal advice, and referral to specialist services. A standard element of good mental health management in construction. [Ch.25]

\item[F10 (CDM Notification)] The form (or online submission) by which a client notifies the HSE of a construction project that exceeds the notification thresholds under CDM 2015 (more than 30 working days with more than 20 workers simultaneously, or more than 500 person-days). The notification is the client's duty. [Ch.27]

\item[FOPS (Falling Object Protective Structure)] A structural protection system fitted to plant cabs (excavators, loaders, dumpers) to protect the operator from objects falling from above. Required where there is a risk of falling objects; FOPS certification indicates compliance with BS EN 13510 or equivalent. [Ch.9]

\item[FORS (Fleet Operator Recognition Scheme)] A voluntary accreditation scheme for vehicle operators setting standards for safety, efficiency, and environment. Increasingly required by principal contractors for logistics providers. [Ch.19, Ch.28]

\item[Genny (Signal Generator)] An electronic instrument used in conjunction with a Cable Avoidance Tool (CAT) to actively inject a signal onto buried services, improving detection reliability. The Genny is connected to a buried service (where accessible) or to a tracer wire and transmits a signal that the CAT can detect. [Ch.6, Ch.22]

\item[H\&S File (Health and Safety File)] A document prepared by the Principal Designer and updated by the Principal Contractor, containing information about the completed building needed for future maintenance and construction work. Must be handed to the client at completion. Required by CDM 2015. [Ch.27]

\item[HAVS (Hand-Arm Vibration Syndrome)] A debilitating and irreversible condition affecting the blood vessels, nerves, and musculoskeletal system of the hand and arm, caused by prolonged exposure to vibrating tools and equipment. Regulated by the Control of Vibration at Work Regulations 2005. [Ch.10]

\item[HAVS Tier Assessment] A structured medical assessment used to diagnose and stage HAVS, typically comprising a standardised questionnaire, clinical examination, and vascular and neurological tests. HAVS tier assessments are conducted by occupational health practitioners. [Ch.10]

\item[Hazard] Anything with the potential to cause harm. In construction, hazards include physical conditions (an unsupported trench, an unprotected edge), energy sources (electricity, pressure, kinetic energy of falling objects), chemical agents (silica dust, solvents, asbestos), biological agents, and organisational/psychosocial conditions (excessive workload, bullying). [Ch.3]

\item[Hierarchy of Controls] The priority order in which risk controls should be considered: Elimination; Substitution; Engineering controls; Administrative controls; PPE. Controls at the top of the hierarchy provide greater and more reliable risk reduction than controls at the bottom. [Ch.4]

\item[Hot Work] Any work involving a source of ignition, including welding (MMA, MIG, TIG), oxy-acetylene cutting, grinding, disc cutting, and the use of blow torches. Hot work requires a Permit to Work and specific fire precautions in any location where there is a risk of fire or explosion. [Throughout]

\item[HSE (Health and Safety Executive)] The UK national regulator and enforcer for workplace health and safety. The HSE has powers of investigation, improvement notice, prohibition notice, and prosecution. The HSE is the primary enforcing authority for CDM-notifiable construction projects. [Throughout]

\item[Improvement Notice] An enforcement notice issued by an HSE or Local Authority inspector requiring the recipient to remedy a specified contravention of health and safety law within a specified time period (minimum 21 days). Failure to comply with an Improvement Notice is a criminal offence. [Ch.2]

\item[Inherent Risk] The level of risk presented by a hazard before any controls are applied. Assessed as Likelihood $\times$ Consequence on the risk matrix. [Ch.3]

\item[IPAF (International Powered Access Federation)] The industry body for the powered access (MEWP) sector, administering the PAL card --- the standard evidence of MEWP operator training and competence. [Ch.7, Ch.15]

\item[Just Culture] An organisational culture in which workers are encouraged to report errors and near misses without fear of automatic punishment, while clear accountability is maintained for gross negligence, reckless behaviour, and deliberate unsafe acts. A prerequisite for effective near-miss reporting. [Ch.30]

\item[LAE (Local Authority Enforcement)] Enforcement of health and safety legislation by Local Authority Environmental Health Officers, as distinct from HSE enforcement. LAEs are responsible for fixed premises and non-CDM-notifiable construction in their jurisdiction. [Ch.2]

\item[Likelihood] The probability that a hazard will cause harm. In the 5$\times$5 risk matrix used in this book, likelihood is rated 1 (highly unlikely) to 5 (almost certain). [Ch.3]

\item[LOLER (Lifting Operations and Lifting Equipment Regulations 1998)] Regulations requiring that lifting equipment used at work is strong and stable enough for the intended load, marked to indicate safe working loads, positioned and installed to minimise risks, used safely, and subject to thorough examination by a competent person at specified intervals. [Ch.15]

\item[LOLER Thorough Examination] The periodic in-depth inspection and testing of lifting equipment and lifting accessories required by LOLER 1998. Frequency depends on equipment type and use; typically every six or twelve months for cranes and accessories, and more frequently for equipment in severe service. [Ch.15]

\item[LOTO (Lock-Out Tag-Out)] A safety procedure for isolating energy sources (electrical, mechanical, hydraulic, pneumatic, thermal, chemical, gravitational) during maintenance or servicing work on machinery and equipment. LOTO requires that the energy source is isolated, locked in the off position, and tagged before work begins. [Ch.13]

\item[MEWP (Mobile Elevating Work Platform)] A self-propelled or towed work platform used to provide temporary access at height. Types include scissor lifts, boom lifts (articulated and telescopic), and vertical personnel lifts. MEWP operators must hold a current IPAF PAL card for the category of MEWP being operated. [Ch.7]

\item[Method Statement] A written document describing the sequence of operations for a specific task, identifying the hazards, the controls, and the specific methods to be used. Combined with the risk assessment, the method statement forms the RAMS. [Throughout]

\item[Near Miss] An event that did not result in injury, illness, or damage but had the potential to do so. Near-miss reporting is a critical element of a proactive safety management system; a high near-miss reporting rate indicates a healthy reporting culture. [Ch.29]

\item[NIHL (Noise-Induced Hearing Loss)] Permanent, irreversible hearing loss caused by prolonged or repeated exposure to high noise levels. The second most common occupational disease in the UK construction industry after HAVS. Regulated by the Control of Noise at Work Regulations 2005. [Ch.10]

\item[PASMA (Prefabricated Aluminium Scaffolding Manufacturers Association)] The industry body for the mobile tower scaffold sector, administering the PASMA card --- the standard evidence of mobile tower scaffold erection and use competence. [Ch.8]

\item[PCI (Pre-Construction Information)] Information gathered by or for the client and provided to designers and contractors before the start of design and construction work. Must include utility records, geotechnical information, information about existing structures, and information about hazardous materials. Required by CDM 2015. [Ch.27]

\item[Permit to Work] A formal documented system for controlling high-risk activities, ensuring that the hazards are identified, the controls are in place, and a competent person has authorised work to proceed. Common permit types in construction include Hot Work, Confined Space, Electrical Isolation, and Work at Height. [Throughout]

\item[PFP (Passive Fire Protection)] Products and systems that protect the building structure and contain fire without active operation, including fire-resisting compartment walls, floors, and doors; cavity barriers; and intumescent coatings on structural steelwork. [Ch.14]

\item[PPE (Personal Protective Equipment)] Equipment worn by a worker to minimise exposure to hazards that cannot be adequately controlled by higher-order controls. PPE is the last resort in the hierarchy of controls. Regulated by the PPE at Work Regulations 2022. [Ch.23]

\item[Principal Contractor (CDM)] The contractor appointed by the client to plan, manage, monitor, and coordinate health and safety during the construction phase on a project involving more than one contractor. The principal contractor must prepare and develop the CPP, ensure cooperation between contractors, and liaise with the Principal Designer. [Ch.27]

\item[Principal Designer (CDM)] The organisation or individual appointed by the client to plan, manage, monitor, and coordinate health and safety during the pre-construction phase. Must be an organisation or individual with control over the pre-construction phase. [Ch.27]

\item[Prohibition Notice] An enforcement notice issued by an HSE or Local Authority inspector requiring the immediate cessation of an activity that involves a risk of serious personal injury. Work must stop immediately on receipt of a Prohibition Notice; it may only be resumed when the notice is withdrawn or lifted on appeal. [Ch.2]

\item[PSSR (Pressure Systems Safety Regulations 2000)] Regulations governing the safe design, installation, use, and maintenance of pressure systems and transportable pressure vessels. Relevant to construction sites where compressed air systems, pressure vessels (autoclaves, boilers, hydraulic systems) are in use. [Ch.9]

\item[PUWER (Provision and Use of Work Equipment Regulations 1998)] Regulations requiring that work equipment is suitable for its intended use, maintained in safe working order, used only by persons trained and competent to use it, and inspected at appropriate intervals. Applies to all work equipment on a construction site. [Ch.9]

\item[RAMS (Risk Assessment and Method Statement)] The combined document package comprising the risk assessment (identifying hazards and evaluating risks) and the method statement (describing the safe method of work). RAMS must be submitted to the principal contractor for review and approval before any new work activity commences. [Throughout]

\item[RCS (Respirable Crystalline Silica)] Fine silica dust generated by cutting, grinding, drilling, or polishing stone, concrete, mortar, brick, or sandstone. Inhalation of RCS can cause silicosis (an irreversible, progressive, and potentially fatal lung disease) and lung cancer. RCS is the highest-priority occupational health dust risk in construction. [Ch.11, Ch.21]

\item[Residual Risk] The level of risk remaining after controls have been applied. The risk assessment must demonstrate that the residual risk is As Low As Reasonably Practicable (ALARP). [Ch.3]

\item[RIDDOR 2013 (Reporting of Injuries, Diseases and Dangerous Occurrences Regulations 2013)] Regulations requiring the reporting of workplace deaths, specified injuries, over-seven-day injuries, dangerous occurrences, and occupational diseases to the relevant enforcing authority. [Ch.29]

\item[Risk Matrix] A tool for evaluating risk by combining Likelihood and Consequence scores. The 5$\times$5 matrix used in this book produces a risk score of 1--25, rated Low (1--4), Medium (5--9), High (10--14), and Very High (15--25). [Ch.3]

\item[ROPS (Roll-Over Protective Structure)] A structural cab or frame fitted to plant (dumpers, telehandlers, tractors) to protect the operator in the event of the machine rolling over. Required by PUWER where roll-over is a foreseeable risk; ROPS certification indicates compliance with BS EN 13510 or equivalent. [Ch.9]

\item[Safe Digging Zone] A zone established either side of an identified or estimated buried service position within which mechanical excavation is prohibited; hand digging only is permitted within the zone. HSG47 recommends 500\,mm either side of the estimated service position. [Ch.6, Ch.22]

\item[SDS (Safety Data Sheet)] A document provided by the manufacturer or supplier of a hazardous chemical product, containing information about the product's hazardous properties, safe handling requirements, PPE requirements, first aid measures, and emergency procedures. Required by COSHH Regulations; must be available at the point of use. [Ch.11]

\item[Specified Injury] A category of injury defined in Schedule 1 to RIDDOR 2013 that must be reported to the relevant enforcing authority. Specified injuries include: fractures (other than fingers, thumbs, and toes); amputations; injuries resulting in loss of consciousness; injuries requiring more than 24 hours in hospital; and others. [Ch.29]

\item[SSIP (Safety Schemes in Procurement)] An umbrella body whose member schemes (CHAS, SafeContractor, Constructionline Gold, Acclaim, and others) assess contractors against a common core criteria set for health and safety competence. SSIP mutual recognition allows assessments conducted by one member scheme to be accepted by another. [Ch.28]

\item[SWL (Safe Working Load)] The maximum load that a piece of equipment, tackle, or structure is rated to carry safely. Often used interchangeably with WLL (Working Load Limit). The SWL must be marked on all lifting equipment and accessories. [Ch.15]

\item[Toolbox Talk] A short, focused safety briefing delivered to a group of operatives at or near the workplace, typically covering a specific hazard, control measure, or recent incident. Toolbox talks are a primary method of communicating health and safety information to the workforce and must be documented with an attendance record. [Throughout]

\item[Vacuum Excavation (Hydrovac / Suction Excavation)] A method of excavation that uses high-pressure water or air to break up the ground and a vacuum system to remove the spoil, enabling safe excavation around buried services without mechanical contact risk. Required by HSG47 where buried service positions are confirmed. [Ch.6, Ch.22]

\item[WLL (Working Load Limit)] The maximum mass that a piece of lifting equipment or tackle is rated to lift safely under specified conditions. See SWL. [Ch.15]

\end{description}
